\scanpage{25}

Le lemme 4 nous ramène alors à vérifier le

\uline{Lemme 5} (\unsure{à énoncer|évidemment} avant le th).

\tikz[remember picture,overlay]{\coordinate (lgfa);}Sous les conditions du th, mais sans \deleted{\unsure{savoir ni a priori que|savoir a priori que}} \inserted{supposer $\deleted{\unsure{x}}y_0$ \deleted{géom.} \uline{unibranche}} en $y_0$, \inserted{ni $f$ propre,} pour tt pt fermé $x$ de $X_0$, \deleted{l'anneau local} $\operatorname{Spec}(\underline{\mathcal O}_{X_0,x})$ est connexe en dim $\geq k$\tikz[remember picture,overlay]{\coordinate (lgfb); \draw ([xshift=-5pt,yshift=4pt]lgfa) -- ([xshift=-5pt]lgfa |- lgfb);}
% slanted
\margin{dire plutôt : tout $\bar x$ tel que $\dim\bar x\leq k$, $\operatorname{Spec}\underline{\mathcal O}_{X_0,x})'$ est connexe}

En effet, le lemme 1 \inserted{+ \unsure{théor.|théorie} de dim} nous ramène facilement, utilisant l'hypothèse sur $X_1$, au cas où $X_1$ est \uline{irréductible} \deleted{de dim $\unsure{>|\geq} k+1$}, donc $X$ \uline{irréductible} dominant $Y$, \inserted{donc $=$ intègre comme $Y$.} On peut supposer par une réduction facile (?) au cas où ``\deleted{la formule des \inserted{\deleted{\unsure{en dépendant de type fini|essentielt de type fini}}} dimensions est vraie pour $\underline{\mathcal O}_{Y,y_0}$}\inserted{$\underline{\mathcal O}_{Y,y_0}$ est formellt équidimensionnel}''.
\margin{expliciter}

Donc

\[
\dim\underline{\mathcal O}_{X,x}=\underbrace{\dim\underline{\mathcal O}_{Y,y_0}}_{n}+\underbrace{\dim X_1}_{\geq k+1}
\]

% archive sheet number: 24
